\chapter{Persamaan Diferensial dan Sistem Dinamik}

Jam bandul di ruang tamu berayun dengan amplitudo enam puluh derajat, dan panjang
lengannya satu meter. Buku pelajaran memberi periode
\[
T=2\pi\sqrt{L/g}=2.006409\ \text{detik},
\]
dan rumus itu memang benar untuk ayunan yang sangat kecil. Diukur pada amplitudo
enam puluh derajat, bandul yang sama memerlukan $2{,}153242$ detik untuk satu ayunan
penuh: $7{,}32$ persen lebih lama. Jam yang disetel memakai rumus kecil itu
ketinggalan $5892$ detik, sekitar sembilan puluh delapan menit, setiap hari.

Rumus itu bukan salah hitung. Ia jawaban tepat untuk persamaan yang berbeda.
Ayunan sebenarnya memenuhi
\[
\frac{d^{2}\theta}{dt^{2}}=-\frac{g}{L}\sin\theta ,
\]
sedangkan rumus periode kecil berasal dari penggantian $\sin\theta$ dengan
$\theta$. Pada $\theta=60^{\circ}$ penggantian itu mengubah $0{,}866025$ menjadi
$1{,}047198$, selisih tujuh belas persen di titik terjauh, dan kesalahan periode
yang tertinggal $7{,}32$ persen. 

Bab 12 mengajari kita menurunkan, Bab 14 deret fungsi, Bab 15 medan vektor.
Tidak satu pun menjawab pertanyaan di sini: bila laju perubahan diketahui,
bagaimana mencari fungsinya, apakah ia tunggal, dan apa yang terjadi dalam jangka
panjang. Bab ini mengisi keempat kekosongan itu: cara menyelesaikan, syarat
keberadaan dan ketunggalan, bidang fase, dan uji kestabilan.

\begin{untukanak}
Bayangkan bak air di kamar mandi. Keran di atas mengalir tetap, lubang pembuangan
di bawah mengalir keluar, dan semakin tinggi air, semakin deras yang keluar.
Kalau kita membuka keran lalu pergi, tinggi air besok tidak dapat dihitung hanya
dari besar keran: yang menentukan adalah tinggi air sekarang. Air yang tinggi
mendorong keluar lebih cepat, kenaikannya melambat, dan akhirnya berhenti ketika
air masuk sama dengan air keluar.

Benda lain bekerja dengan cara yang sama. Kopi panas mendingin lebih cepat ketika
masih sangat panas. Kuman berkembang biak lebih cepat ketika jumlahnya banyak,
lalu melambat ketika makanan habis. Yang kita tahu adalah aturan perubahan, bukan
hasil akhirnya.

Sekarang pertanyaan yang membuat pekerjaan ini sulit. Dua bak sama, keran sama,
lubang sama, tetapi satu mulai kosong dan satu mulai setengah penuh: apakah
akhirnya tingginya sama? Kadang ya, kadang tidak. Dan ada lubang yang jahat,
tersumbat sebagian lalu terbuka tiba-tiba: dua cerita berbeda bisa keluar dari
keadaan awal yang sama.

Di bab ini kita akan tahu mengapa aturan perubahan kadang cukup menentukan
seluruh masa depan, dan kadang tidak.
\end{untukanak}

\section{Laju yang bergantung pada dirinya sendiri}
\tingkat{dasar}

Persamaan diferensial menyamakan laju perubahan sebuah fungsi dengan sesuatu
yang memuat fungsi itu sendiri. Menyelesaikannya berarti mencari fungsi yang
memenuhi kalimat itu pada seluruh sebuah selang. Turunan di Bab 12 dihitung di
satu titik; solusi di sini adalah fungsi pada selang, dan perbedaan itu yang
paling sering terlewat.

\begin{definisi}[Persamaan diferensial biasa]\index{persamaan diferensial}
Persamaan
\[
F\bigl(t,y,y',y'',\dots,y^{(n)}\bigr)=0
\]
disebut \emph{persamaan diferensial biasa} (PDB) orde $n$ untuk fungsi $y$ dari
peubah $t$ bila $F$ memuat turunan ke-$n$ dan tidak memuat turunan yang lebih
tinggi. PDB itu disebut \emph{linear} bila $F$ berbentuk
\[
a_{n}(t)\,y^{(n)}+\cdots+a_{1}(t)\,y'+a_{0}(t)\,y=b(t),
\]
dan disebut \emph{otonom} bila $F$ tidak memuat $t$ secara eksplisit.
\end{definisi}

 Semua contoh di
bab ini orde satu atau orde dua, karena keduanya sudah memuat seluruh gagasan
penting.

\begin{definisi}[Solusi dan masalah nilai awal]\index{masalah nilai awal}
Fungsi $y$ disebut \emph{solusi} PDB pada selang $I$ bila $y$ mempunyai turunan
sampai orde $n$ pada $I$ dan memenuhi persamaan di setiap titik $I$. Bila
ditambahkan syarat $y(t_{0})=y_{0}$, gabungan keduanya disebut \emph{masalah
nilai awal}, dan $y_{0}$ disebut nilai awal.
\end{definisi}

Satu contoh untuk memastikan definisinya bekerja. Ambil model populasi
\[
y'=y(1-y),\qquad y(0)=\tfrac12 ,
\]
dengan $y$ menyatakan bagian populasi terhadap daya tampung. Calon solusinya
\[
y(t)=\frac{1}{1+e^{-t}} .
\]
Turunannya $y'(t)=e^{-t}/(1+e^{-t})^{2}$, dan ruas kanan
$y(1-y)=e^{-t}/(1+e^{-t})^{2}$ sama persis. Nilai awalnya $y(0)=1/2$ juga cocok.
Jadi fungsinya solusi, dan pemeriksaan itu hanya memakai aturan turunan hasil
bagi dari Bab 12.

Yang perlu diperhatikan: $y(t)=1$ dan $y(t)=0$ juga solusi, keduanya konstan,
dan keduanya muncul saat kita menanyakan kapan ruas kanan nol. Jadi satu
persamaan dapat punya banyak solusi; nilai awal memilih satu di antaranya.
Koleksi semua solusi disebut \emph{solusi umum}, dan ia memuat satu konstanta
untuk setiap orde.

Selang tempat solusi hidup adalah bagian dari jawaban. Ambil
\[
y'=y^{2},\qquad y(0)=1 .
\]
Pemisahan variabel memberi $y(t)=1/(1-t)$, dan rumus itu hanya masuk akal pada
$t<1$. Di $t=0{,}9$ nilainya $10$; di $t=0{,}99$ nilainya $100$; di $t=0{,}999$ nilainya
$1000$. Tidak ada bilangan berhingga di $t=1$, dan tidak ada solusi yang
melanjutkan fungsi itu melewatinya. Solusinya \emph{meledak} dalam waktu
berhingga, dan program yang melaporkan nilai di $t=2$ tidak menemukan solusi; ia
menemukan pembulatan.

\kamus{persamaan diferensial}{differential equation}{Kalimat yang mengikat laju perubahan sebuah besaran dengan besar besaran itu sendiri.}

\kamus{solusi umum}{general solution}{Keluarga semua fungsi yang memenuhi persamaan, biasanya memuat satu konstanta bebas per orde.}

\begin{bacakeras}
--- Kalau bola ini kulepas di lereng, ia akan bergerak ke bawah. Berapa cepat?
--- Tergantung di mana bola itu sekarang. Di dekat puncak, pelan. Di tengah
lereng, cepat.
--- Jadi kecepatannya bukan satu angka tetap?
--- Bukan. Ia aturan: di tiap titik, kecepatan ditentukan oleh titik itu.
--- Kalau aturannya sudah tahu segalanya, apakah jalannya sudah tertentu?
--- Belum. Aturan memberi arah di setiap titik; jalan lahir dari titik awal yang
kau pilih.
--- Dua bola, dua titik awal, satu aturan?
--- Satu aturan, dua jalan.
--- Kalau dua jalan bertemu di titik yang sama, bagaimana kita tahu mana yang
datang dari mana?
--- Kita tidak bisa. Karena itu bab ini menuntut satu titik awal untuk satu
jalan.

Ulangi sekali, dengan suara biasa: aturan sama, titik awal berbeda, jalan berbeda.
\end{bacakeras}

\section{Dua keluarga yang dapat diselesaikan dengan tangan}
\tingkat{dasar}

Tidak semua persamaan diferensial punya rumus penyelesaian, tetapi dua keluarga
besar punya, dan keduanya muncul di hampir setiap penerapan tingkat pertama.
Keluarga pertama memisahkan peubah; keluarga kedua linear orde satu.

\subsection{Pemisahan variabel}

Bila persamaan dapat ditulis $y'=f(t)g(y)$, kita memindahkan semua $y$ ke satu
sisi dan semua $t$ ke sisi lain:
\[
\frac{dy}{g(y)}=f(t)\,dt,\qquad \int\frac{dy}{g(y)}=\int f(t)\,dt .
\]
Langkah itu bukan trik notasi, melainkan aturan rantai yang dibalik: bila $H$
antiturunan $1/g$, maka $\frac{d}{dt}H(y(t))=H'(y)y'=y'/g(y)$, sehingga
persamaan asal menjadi $\frac{d}{dt}H(y(t))=f(t)$ dan Bab 13 memberi
$H(y)=\int f\,dt+C$. Hasilnya biasanya implisit, dan itu jawaban yang sah.

Model logistik di atas dipakai sepanjang bab. Dengan
$g(y)=y(1-y)$ dan pecahan parsial,
\[
\int\frac{dy}{y(1-y)}=\int\Bigl(\frac{1}{y}+\frac{1}{1-y}\Bigr)dy
=\ln\Bigl|\frac{y}{1-y}\Bigr| ,
\]
maka $\ln|y/(1-y)|=t+C$ dan
\[
y(t)=\frac{1}{1+Ae^{-t}},\qquad A=\frac{1-y_{0}}{y_{0}} .
\]
Untuk $y_{0}=1/2$ diperoleh $A=1$, yaitu rumus yang kita periksa tadi. Perhatikan
bahwa $A$ ditentukan oleh nilai awal, dan bahwa $A=0$ memberi solusi konstan $y=1$
yang tidak tercakup pembagian di atas: saat memisahkan variabel kita membagi
dengan $g(y_{0})$, dan pembagian itu hanya sah bila $g(y_{0})\ne0$.

\subsection{Persamaan linear dan faktor integrasi}

Bentuk baku persamaan linear orde satu adalah
\[
y'+p(t)\,y=q(t).
\]
Ruas kiri hampir merupakan turunan sebuah hasil kali, hanya kurang satu faktor.
Faktor yang hilang itulah yang dicari teorema berikut.

\begin{teorema}[Faktor integrasi]\index{faktor integrasi}
Misalkan $p$ dan $q$ kontinu pada selang buka $I$ yang memuat $t_{0}$, dan
misalkan
\[
\mu(t)=\exp\Bigl(\int_{t_{0}}^{t}p(s)\,ds\Bigr).
\]
Maka masalah nilai awal $y'+p(t)y=q(t)$, $y(t_{0})=y_{0}$ mempunyai tepat satu
solusi pada $I$, yaitu
\[
y(t)=\frac{1}{\mu(t)}\Bigl(y_{0}+\int_{t_{0}}^{t}\mu(s)\,q(s)\,ds\Bigr).
\]
\end{teorema}

\begin{proof}
Karena $p$ kontinu, $\mu$ terdefinisi, positif, dan terdiferensialkan pada $I$,
dengan $\mu'=p\mu$ menurut teorema dasar kalkulus (Bab 13). Kalikan persamaan
dengan $\mu$:
\[
\mu y'+\mu p\,y=\mu q
\quad\Longleftrightarrow\quad
(\mu y)'=\mu q ,
\]
karena aturan turunan hasil kali memberi $(\mu y)'=\mu'y+\mu y'=\mu(py+y')$.
Integralkan kedua ruas dari $t_{0}$ sampai $t$ dan pakai teorema dasar kalkulus:
\[
\mu(t)y(t)-\mu(t_{0})y(t_{0})=\int_{t_{0}}^{t}\mu(s)q(s)\,ds .
\]
Karena $\mu(t_{0})=1$, pembagian dengan $\mu(t)>0$ memberi rumus yang diminta.
Ketunggalan: bila $y_{1},y_{2}$ dua solusi, selisihnya $w=y_{1}-y_{2}$ memenuhi
$w'+pw=0$ dengan $w(t_{0})=0$; mengalikan dengan $\mu$ memberi $(\mu w)'=0$,
jadi $\mu w$ konstan, dan karena nilainya nol di $t_{0}$ ia nol di mana-mana.
\end{proof}

Bukti itu tidak menebak apa pun: $\mu$ dicari dengan menuntut $\mu'=p\mu$,
persamaan yang lebih mudah daripada asalnya. Pola ini berulang: satu persoalan
diselesaikan dengan mengubahnya menjadi persoalan yang lebih sederhana.

\begin{contoh}[Satu persamaan linear, dua pemeriksaan]
Selesaikan $y'+2y=4t$ dengan $y(0)=1$. Di sini $p=2$ dan $q=4t$, sehingga
$\mu=e^{2t}$ dan
\[
(e^{2t}y)'=4t\,e^{2t}.
\]
Integral parsial memberi
\[
\int 4t\,e^{2t}\,dt=\Bigl(2t-1\Bigr)e^{2t}+C ,
\]
sebab turunan $(2t-1)e^{2t}$ adalah $2e^{2t}+2(2t-1)e^{2t}=(4t)e^{2t}$. Jadi
\[
y(t)=2t-1+Ce^{-2t},
\]
dan syarat $y(0)=1$ memberi $C=2$, yaitu $y(t)=2t-1+2e^{-2t}$.

Pemeriksaan: substitusi memberi $y'+2y=4t$, dan $y(1)=1+2e^{-2}=1{,}2706705665$;
integrator adaptif dengan toleransi $10^{-12}$ memberi $1{,}270670566473$.
\end{contoh}

\subsection{Rangkaian RC: contoh nyata yang paling sederhana}

Sebuah sumber tegangan $V_{s}$ dihubungkan seri dengan resistor $R$ dan
kapasitor $C$. Hukum Kirchhoff tentang tegangan memberi $V_{s}=Ri+V_{C}$, dan
hukum kapasitor memberi $i=C\,dV_{C}/dt$. Keduanya memberi
\[
RC\,\frac{dV_{C}}{dt}+V_{C}=V_{s},
\qquad\text{atau}\qquad
\frac{dV_{C}}{dt}+\frac{1}{RC}V_{C}=\frac{V_{s}}{RC}.
\]
Ini persamaan linear dengan $p=1/(RC)$ konstan. Menurut teorema tadi, dengan
$V_{C}(0)=0$,
\[
V_{C}(t)=V_{s}\bigl(1-e^{-t/\tau}\bigr),\qquad \tau=RC .
\]
Angka $\tau$ disebut tetapan waktu. Untuk $R=1\,\mathrm{k}\Omega$ dan
$C=1\,\mu\mathrm{F}$, $\tau=10^{-3}$ detik. Dengan $V_{s}=5$ volt, tegangan
pada $t=\tau$ adalah $3{,}160603$ volt dan pada $t=5\tau$ adalah $4{,}966310$ volt,
yaitu $99{,}33$ persen dari $5$ volt. Untuk mencapai $90$ persen dibutuhkan
$t=\tau\ln10=2{,}302585$ milidetik.

\begin{resep}
\begin{enumerate}
\item Tulis persamaan dalam bentuk baku $y'+p(t)y=q(t)$: bagi dengan koefisien
  $y'$ sehingga koefisiennya tepat satu.
\item Hitung $\mu(t)=\exp\bigl(\int p(t)\,dt\bigr)$, tanpa konstanta integrasi.
\item Kalikan seluruh persamaan dengan $\mu$, lalu tulis ruas kiri sebagai
  $(\mu y)'$. Bila ruas kiri tidak mau menjadi satu turunan, langkah 1 salah.
\item Integralkan kedua ruas terhadap $t$, lalu bagi dengan $\mu$.
\item Pakai nilai awal untuk menentukan konstanta, dan tulis selang berlakunya.
\end{enumerate}
Peringatan: langkah 1 tidak boleh dilewati, dan langkah 3 tidak boleh dikerjakan
sebelum langkah 2. Menghitung $\mu$ dari koefisien yang belum dinormalkan memberi
faktor integrasi yang salah, dan kesalahannya baru kelihatan saat solusinya
diuji.
\end{resep}

\kamus{faktor integrasi}{integrating factor}{Pengali yang mengubah ruas kiri persamaan linear menjadi satu turunan hasil kali.}

\section{Keberadaan, ketunggalan, dan iterasi Picard}
\tingkat{menengah}

Contoh $y'=y^{2}$ menunjukkan solusi bisa berhenti ada. Diberi $f$ dan satu titik
awal, apakah solusinya ada dan apakah hanya satu? Jawabannya bergantung pada satu
sifat $f$ terhadap $y$.

\begin{definisi}[Syarat Lipschitz]\index{syarat Lipschitz}
Fungsi $f(t,y)$ memenuhi syarat Lipschitz terhadap $y$ pada daerah
$D\subseteq\mathbb{R}^{2}$ bila ada konstanta $L$ sehingga
\[
|f(t,y_{1})-f(t,y_{2})|\le L\,|y_{1}-y_{2}|
\]
untuk setiap $(t,y_{1}),(t,y_{2})\in D$. Konstanta terkecil yang mungkin disebut
konstanta Lipschitz.
\end{definisi}

Kalimat sehari-harinya: $f$ tidak boleh berubah terlalu mendadak ketika $y$
bergeser sedikit. Bila $f$ punya turunan parsial terhadap $y$ yang kontinu dan
terbatas oleh $L$ pada $D$, syarat itu otomatis berlaku, karena teorema nilai
rata-rata (Bab 12) memberi
$f(t,y_{1})-f(t,y_{2})=\partial_{y}f(t,\xi)(y_{1}-y_{2})$.

\begin{teorema}[Picard--Lindel\"of]\index{teorema Picard-Lindelof}
Misalkan $f$ kontinu pada persegi panjang tertutup
$D=[t_{0}-a,t_{0}+a]\times[y_{0}-b,y_{0}+b]$ dan memenuhi syarat Lipschitz
terhadap $y$ pada $D$ dengan konstanta $L$. Misalkan $M=\max_{D}|f|$ dan
$h=\min(a,\,b/M)$. Maka masalah nilai awal $y'=f(t,y)$, $y(t_{0})=y_{0}$
mempunyai tepat satu solusi pada selang $|t-t_{0}|\le h$.
\end{teorema}

Bukti lengkapnya memakai teorema titik tetap Banach pada ruang fungsi kontinu
dengan norma maksimum, yang lengkap karena kekontinuan seragam dibahas di Bab 11
dan Bab 14. Saya berikan sketsanya, karena sketsa itu memuat seluruh gagasan dan
dapat dijadikan algoritma.

Tulis persamaan sebagai persamaan integral, yang setara karena kedua sisinya
kontinu:
\[
y(t)=y_{0}+\int_{t_{0}}^{t}f\bigl(s,y(s)\bigr)\,ds .
\]
Definisikan pemetaan $T$ dengan rumus itu, mulai dari $y_{0}(t)=y_{0}$ konstan,
lalu ulangi:
\[
y_{n+1}(t)=y_{0}+\int_{t_{0}}^{t}f\bigl(s,y_{n}(s)\bigr)\,ds .
\]
Setiap $y_{n}$ terletak di dalam $D$ selama $|t-t_{0}|\le h$, karena
$|y_{n+1}-y_{0}|\le M|t-t_{0}|\le Mh\le b$. Selisih dua iterasi berturut-turut
memenuhi
\[
|y_{n+1}(t)-y_{n}(t)|\le L\Bigl|\int_{t_{0}}^{t}|y_{n}(s)-y_{n-1}(s)|\,ds\Bigr|
\le \frac{(L|t-t_{0}|)^{n}}{n!}\,\|y_{1}-y_{0}\| ,
\]
dan induksi pada $n$ memberi batas itu. Karena $\sum_{n}(Lh)^{n}/n!$ konvergen
(uji rasio, Bab 5), deret $y_{0}+\sum(y_{n+1}-y_{n})$ konvergen seragam, dan
batasnya adalah solusi. Ketunggalan datang dari pertidaksamaan yang sama: dua
solusi berbeda akan memenuhi $|y_{1}-y_{2}|\le L|t-t_{0}|\,\|y_{1}-y_{2}\|$,
yang untuk $Lh<1$ hanya mungkin bila selisihnya nol; kasus $Lh\ge1$ ditangani
dengan memperkecil $h$, dan itulah sebab teorema hanya menjanjikan selang
pendek.

Jalankan iterasi itu pada $y'=y$, $y(0)=1$. Mulai dari $y_{0}=1$:
\[
y_{1}=1+t,\quad y_{2}=1+t+\frac{t^{2}}{2},\quad
y_{3}=1+t+\frac{t^{2}}{2}+\frac{t^{3}}{6},
\]
sehingga $y_{n}$ adalah jumlah parsial deret eksponensial dari Bab 5. Di $t=1$,
iterasi kedelapan memberi $2{,}718278769841$, sedangkan $e=2{,}718281828459$;
galatnya $3{,}06\times10^{-6}$.

Syarat Lipschitz tidak boleh dibuang. Ambil
\[
y'=\sqrt{|y|},\qquad y(0)=0 .
\]
Ruas kanannya kontinu, jadi solusi ada; tetapi ia tidak memenuhi syarat Lipschitz
di sekitar $y=0$, karena kemiringannya tak terbatas di sana. Dan memang ada dua
solusi: fungsi nol, $y(t)\equiv0$, dan
\[
y(t)=\frac{t^{2}}{4}\quad (t\ge0).
\]
Pemeriksaannya satu baris: $y'=t/2$ dan $\sqrt{|y|}=\sqrt{t^{2}/4}=t/2$ untuk
$t\ge0$, sedangkan $y(0)=0$. Dua riwayat berbeda dari satu keadaan awal yang
sama, persis seperti lubang bak yang terbuka tiba-tiba pada kotak pembuka bab.

Lipschitz adalah syarat \emph{cukup}, bukan perlu. Keberadaan tanpa
ketunggalan dijamin teorema Peano, yang memakai kekompakan (Bab 11) dan
Arzel\`a--Ascoli; rujukannya Coddington dan Levinson, bab 1--2.

\kamus{ketunggalan}{uniqueness}{Jaminan bahwa satu keadaan awal hanya melahirkan satu riwayat, tanpa cabang.}
\section{Sistem, bidang fase, dan osilator}
\tingkat{menengah}

Satu persamaan orde dua bukan binatang baru; ia dua persamaan orde satu yang
berjalan bersama. Dengan menamai $v=\theta'$, persamaan pendulum menjadi
\[
\theta'=v,\qquad v'=-\frac{g}{L}\sin\theta .
\]
Pasangan ini otonom: ruas kanannya tidak memuat $t$, sehingga keadaan awal
menentukan seluruh riwayat (bila syarat Lipschitz dipenuhi) dan kita boleh
melupakan waktu sebentar untuk bertanya apa yang terjadi pada bidang
$(\theta,v)$.

\begin{definisi}[Sistem otonom dan bidang fase]\index{bidang fase}
Sistem
\[
\mathbf{x}'=\mathbf{f}(\mathbf{x}),\qquad \mathbf{x}\in\mathbb{R}^{n},
\]
disebut \emph{otonom}. Untuk $n=2$, bidang tempat titik $\mathbf{x}$ hidup
disebut \emph{bidang fase}; kurva $\{\mathbf{x}(t)\}$ disebut \emph{lintasan},
dan gambar seluruh panah $\mathbf{f}(\mathbf{x})$ disebut \emph{medan arah}.
Titik $\mathbf{x}^{*}$ dengan $\mathbf{f}(\mathbf{x}^{*})=0$ disebut
\emph{titik setimbang} atau titik tetap.
\end{definisi}

Sumbu mendatar bidang fase menyatakan posisi, sumbu tegak menyatakan laju, dan
satu titik berarti "di sini, dengan kecepatan ini". Waktu tidak digambar; ia
hanya parameter yang menyusuri kurva. Itulah keuntungannya: seluruh riwayat, untuk
semua titik awal sekaligus, tampak sebagai satu gambar.

Titik setimbang pendulum mudah dicari. $\theta'=0$ menuntut $v=0$, dan $v'=0$
menuntut $\sin\theta=0$, jadi $\theta=k\pi$. Di $\theta=0$ bandul tergantung
lurus ke bawah; di $\theta=\pi$ ia berdiri tegak ke atas. Yang pertama stabil
dalam pengertian sehari-hari, yang kedua tidak: sentuhan sekecil apa pun
menjatuhkannya.

\begin{untukanak}
Pergi ke taman, cari ayunan, dan duduklah. Minta seorang teman mendorongmu
pelan, hanya sedikit. Hitung sepuluh ayunan dengan menghitung di dalam kepala:
satu, dua, tiga, sampai sepuluh. Sekarang minta dorongan yang jauh lebih tinggi,
sampai ayunan naik hampir setinggi mistar atas, lalu hitung sepuluh ayunan lagi.

Yang akan kau lihat: ayunan tinggi memerlukan waktu lebih lama untuk satu
putaran penuh, walaupun kau merasa seolah-olah ayunan tinggi lebih cepat. Rasa
itu menipu karena jarak yang ditempuh jauh lebih panjang, sedangkan waktu satu
putaran hampir sama.

Sekarang lihat bagian atas ayunan. Saat kau berada di titik tertinggi, ada satu
saat pendek ketika kau hampir berhenti dan menggantung. Di titik terendah, kau
melaju paling cepat. Ayunan tidak pernah berhenti di tengah, dan tidak pernah
paling cepat di ujung. Dua hal itu selalu bertukar tempat.

Di bab ini kita akan tahu mengapa ayunan tinggi lebih lambat, dan mengapa
bandul yang berdiri tegak di atas hampir tidak mungkin dipertahankan.
\end{untukanak}

Sekarang rumus periode yang tepat. Kalikan $v'=-(g/L)\sin\theta$ dengan $v=\theta'$
dan integrasikan:
\[
\frac{d}{dt}\Bigl(\frac{v^{2}}{2}+\frac{g}{L}(1-\cos\theta)\Bigr)=0 .
\]
Kurung itu adalah energi per satuan massa, dan ia konstan sepanjang lintasan.
Dengan $E$ menyatakan nilainya,
\[
\frac{v^{2}}{2}+\frac{g}{L}\bigl(1-\cos\theta\bigr)=E ,
\]
sehingga $v=\pm\sqrt{2E-2(g/L)(1-\cos\theta)}$. Lintasan di bidang fase adalah
kurva-kurva energi itu; pada Gambar~\ref{fig:bidang-fase-pendulum} mereka
digambar untuk satuan $g/L=1$.

\begin{figure}[htbp]
\centering
\begin{tikzpicture}[x=1.05cm,y=1.0cm,font=\footnotesize]
  \draw[->,layergrey] (-4.15,0) -- (4.15,0) node[right] {$\theta$};
  \draw[->,layergrey] (0,-2.95) -- (0,2.95) node[above] {$v$};
  \foreach \tick/\lab in {-3.1416/{$-\pi$},-1.5708/{$-\pi/2$},
                          1.5708/{$\pi/2$},3.1416/{$\pi$}}{
    \draw[layergrey!55] (\tick,0.07) -- (\tick,-0.07);
    \node[layergrey,font=\scriptsize,below] at (\tick,-0.07) {\lab};}
  \node[layergrey,font=\scriptsize,below left] at (-0.06,0.02) {$0$};
  \foreach \th in {-3.5,-2.6,-1.7,-0.85,0.85,1.7,2.6,3.5}{
    \foreach \w in {-2.4,-1.6,-0.8,0,0.8,1.6,2.4}{
      \pgfmathsetmacro{\nn}{sqrt(\w*\w+sin(\th r)*sin(\th r))}
      \pgfmathsetmacro{\dx}{0.32*\w/\nn}
      \pgfmathsetmacro{\dy}{-0.32*sin(\th r)/\nn}
      \draw[-{Latex[length=0.9mm]},layerteal!60]
        (\th,\w) -- ({\th+\dx},{\w+\dy});}}
  \foreach \En/\tm in {0.4/0.9273,1.0/1.5708,1.6/2.2143}{
    \draw[layerindigo,thick]
      plot[domain=-\tm:\tm,samples=121,variable=\t]
        ({\t},{sqrt(max(0,2*(\En-1+cos(\t r))))});
    \draw[layerindigo,thick]
      plot[domain=-\tm:\tm,samples=121,variable=\t]
        ({\t},{-sqrt(max(0,2*(\En-1+cos(\t r))))});}
  \draw[accent,very thick]
    plot[domain=-3.1416:3.1416,samples=161,variable=\t]
      ({\t},{2*cos(\t r/2)});
  \draw[accent,very thick]
    plot[domain=-3.1416:3.1416,samples=161,variable=\t]
      ({\t},{-2*cos(\t r/2)});
  \draw[layerplum,thick,densely dashed]
    plot[domain=-3.6:3.6,samples=161,variable=\t]
      ({\t},{sqrt(2*(1.6+cos(\t r)))});
  \fill[ink] (0,0) circle (1.6pt);
  \node[ink,anchor=north east,font=\scriptsize] at (-0.12,-0.52) {pusat};
  \fill[accent] (3.1416,0) circle (1.6pt);
  \fill[accent] (-3.1416,0) circle (1.6pt);
  \node[accent,anchor=north west,font=\scriptsize,align=left]
    at (2.35,-0.18) {pelana};
  \node[accent,anchor=west,font=\scriptsize,align=left]
    at (2.30,1.30) {separatriks $E=2$};
  \node[layerplum,anchor=west,font=\scriptsize,align=left]
    at (0.55,2.72) {lintasan terbuka $E=2.6$};
\end{tikzpicture}
\caption{Bidang fase pendulum dalam satuan $g/L=1$, dengan sumbu mendatar sudut
$\theta$ dan sumbu tegak kecepatan $v$. Panah kecil menunjukkan medan arah
$(\theta',v')=(v,-\sin\theta)$. Kurva biru tua adalah tiga lintasan tertutup dengan
energi $E=0{,}4$, $1{,}0$, dan $1{,}6$; makin besar $E$, makin lebar putarannya.
Kurva merah adalah separatriks $E=2$, yaitu lintasan yang tepat mencapai
$\theta=\pm\pi$ dengan kecepatan nol; ia memisahkan ayunan tertutup dari putaran
terbuka. Kurva putus-putus adalah lintasan terbuka $E=2{,}6$, tempat bandul
berputar penuh tanpa pernah berhenti. Titik di $\theta=0$ adalah pusat, titik di
$\theta=\pm\pi$ adalah pelana.}
\label{fig:bidang-fase-pendulum}
\end{figure}

Gambar itu memuat seluruh perilaku kualitatif pendulum tanpa satu persamaan
diselesaikan; yang dibutuhkan hanya tanda $\sin\theta$ dan simetri energi.

Untuk amplitudo kecil, $\sin\theta\approx\theta$ dan sistemnya menjadi
osilator harmonik
\[
\theta'=v,\qquad v'=-\omega_{0}^{2}\theta,\qquad \omega_{0}=\sqrt{g/L},
\]
dengan solusi $\theta(t)=\theta_{0}\cos(\omega_{0}t)$ dan periode
$2\pi/\omega_{0}=2{,}006409$ detik untuk $L=1$ meter. Lintasannya elips
$v^{2}+\omega_{0}^{2}\theta^{2}=\text{konstan}$, dan itulah tiga lingkaran
kecil di dekat pusat pada gambar tadi.

Periode yang tepat diperoleh dari kekekalan energi, bukan dari linearisasi.
Dengan $E=(g/L)(1-\cos\theta_{0})$ untuk bandul yang dilepas dari diam di
$\theta_{0}$, waktu turun dari $\theta_{0}$ ke nol adalah
\[
\frac{T}{4}=\sqrt{\frac{L}{2g}}\int_{0}^{\theta_{0}}
\frac{d\theta}{\sqrt{\cos\theta-\cos\theta_{0}}}
=\sqrt{\frac{L}{g}}\,K(k),\qquad k=\sin\frac{\theta_{0}}{2},
\]
dengan $K$ integral eliptik lengkap jenis pertama. Untuk $\theta_{0}=60^{\circ}$
diperoleh $k=0{,}5$ dan $K$ pada parameter $m=k^{2}=0{,}25$ bernilai $1{,}6857503548$,
sehingga $T=2{,}153242$ detik. Kuadratur pada integral itu memberi angka yang sama
sampai $9\times10^{-13}$; dua jalan berbeda, satu hasil. Untuk $\theta_{0}=5^{\circ}$
rasionya $1{,}00047617$, yaitu kesalahan kurang dari seperdua ribu: di situlah
rumus kecil pantas dipakai, dan di situlah batas keberlakuannya.

\begin{tangan}
Ambil tali sepanjang satu meter, sebuah mur logam sebagai beban, dan stopwatch
di telepon. Ikat mur di ujung tali, tempelkan ujung lain ke kusen pintu dengan
selotip, dan ukur panjang tali dari titik gantung ke pusat mur: usahakan tepat
satu meter. Tarik mur ke samping hanya sekitar lima sentimeter, lepaskan, dan
catat waktu untuk sepuluh ayunan penuh. Bagi sepuluh; hasilnya sekitar dua
detik.

Sekarang tarik jauh lebih tinggi, kira-kira setinggi setengah meter dari titik
terendah, dan ulangi. Bandingkan kedua angka. Yang akan kau lihat: angka kedua
lebih besar, tetapi tidak dua kali lebih besar. Selisihnya sekitar tujuh persen,
dan itu cocok dengan hitungan di atas bila sudutnya sekitar enam puluh derajat.

Ulangi dengan tali seperempat meter: periodenya menjadi setengah, karena
periode bergantung pada akar panjang.
\end{tangan}

\section{Kestabilan: nilai eigen, linearisasi, dan Lyapunov}
\tingkat{lanjut}

Pertanyaan yang paling sering diajukan tentang sebuah sistem bukan "berapa
nilainya pada $t=7$", melainkan "apakah ia kembali setelah diganggu". Untuk
menjawabnya, kita tidak perlu solusi; cukup turunan di sekitar titik setimbang.

\begin{definisi}[Kestabilan]\index{kestabilan}
Misalkan $\mathbf{x}^{*}$ titik setimbang $\mathbf{x}'=\mathbf{f}(\mathbf{x})$.
Titik itu disebut \emph{stabil} bila untuk setiap $\varepsilon>0$ ada
$\delta>0$ sehingga $\|\mathbf{x}(0)-\mathbf{x}^{*}\|<\delta$ menjamin
$\|\mathbf{x}(t)-\mathbf{x}^{*}\|<\varepsilon$ untuk setiap $t\ge0$. Titik itu
disebut \emph{stabil asimtotik} bila stabil dan ada $\delta>0$ sehingga
$\mathbf{x}(t)\to\mathbf{x}^{*}$ untuk $t\to\infty$. Bila tidak stabil, ia
disebut \emph{tidak stabil}.
\end{definisi}

\begin{teorema}[Linearisasi]\index{linearisasi}
Misalkan $\mathbf{f}$ mempunyai turunan parsial kontinu di sekitar titik
setimbang $\mathbf{x}^{*}$ dan misalkan $A=J_{\mathbf{f}}(\mathbf{x}^{*})$
matriks Jacobian dari Bab 12. Bila tidak ada nilai eigen $A$ yang bagian
realnya nol (titik setimbang \emph{hiperbolik}), maka perilaku kualitatif sistem
nonlinear di sekitar $\mathbf{x}^{*}$ sama dengan perilaku sistem linear
$\mathbf{u}'=A\mathbf{u}$: jenis kestabilannya sama, dan lintasan di dekatnya
dapat dipetakan satu-satu dan kontinu ke lintasan linear.
\end{teorema}

Ini teorema Hartman--Grobman; buktinya panjang dan memakai teorema titik tetap
pada ruang fungsi, jadi saya kutip: Hartman, bab 9, Perko, bagian 2.8. Akibat
praktisnya: kestabilan lokal dibaca dari nilai eigen.

Untuk sistem dua dimensi, seluruh kemungkinan dibaca dari jejak
$\tau=\operatorname{tr}A$ dan determinan $\Delta=\det A$; nilai eigennya akar
$\lambda^{2}-\tau\lambda+\Delta=0$. Bila $\Delta<0$, akarnya real dan
berlawanan tanda: \emph{pelana}. Bila $\Delta>0$ dan $\tau^{2}\ge4\Delta$,
keduanya real dan bertanda sama: \emph{simpul}, stabil bila $\tau<0$. Bila
$\tau^{2}<4\Delta$, akarnya konjugat dengan bagian real $\tau/2$:
\emph{spiral}, stabil bila $\tau<0$, dan pusat bila $\tau=0$.

Contoh: matriks $A=\begin{pmatrix}0&1\\-2&-3\end{pmatrix}$ punya $\tau=-3$
dan $\Delta=2$, sehingga akarnya $\lambda^{2}+3\lambda+2=0$, yaitu $-1$ dan
$-2$: simpul stabil. Matriks $B=\begin{pmatrix}0&1\\-2&0\end{pmatrix}$ punya
$\tau=0$ dan $\Delta=2>0$, akarnya $\pm i\sqrt{2}$: pusat.

Gambar~\ref{fig:potret-linear} memperlihatkan dua potret yang paling berbeda
sifatnya: pelana, tempat hampir semua titik awal menjauh, dan spiral stabil,
tempat semuanya berputar masuk.

\begin{figure}[htbp]
\centering
\begin{tikzpicture}[font=\footnotesize]
  \begin{scope}[x=1.35cm,y=1.35cm]
    \begin{scope}
      \clip (-1.62,-1.52) rectangle (1.62,1.52);
      \foreach \c in {0.05,0.30,0.70,1.20}{
        \draw[layerteal,thick] plot[domain=-1.55:1.55,samples=90,variable=\s]
          ({sqrt(\c+\s*\s)},{\s});
        \draw[layerteal,thick] plot[domain=-1.55:1.55,samples=90,variable=\s]
          ({-sqrt(\c+\s*\s)},{\s});}
      \draw[layerindigo,densely dashed] (-1.5,-1.5) -- (1.5,1.5);
      \draw[layerindigo,densely dashed] (-1.5,1.5) -- (1.5,-1.5);
    \end{scope}
    \draw[->,layergrey] (-1.75,0) -- (1.75,0) node[right] {$x$};
    \draw[->,layergrey] (0,-1.55) -- (0,1.55) node[above] {$y$};
    \fill[ink] (0,0) circle (1.5pt);
    \node[layerindigo,font=\scriptsize,anchor=south east] at (-1.52,1.30)
      {arah eigen};
    \node[layergrey,anchor=north] at (0,-1.62)
      {(a) pelana: $\lambda=1,-1$};
  \end{scope}
  \begin{scope}[xshift=6.6cm,x=1.35cm,y=1.35cm]
    \draw[->,layergrey] (-1.75,0) -- (1.75,0) node[right] {$x$};
    \draw[->,layergrey] (0,-1.55) -- (0,1.55) node[above] {$y$};
    \foreach \r in {0.35,0.7,1.1,1.5}{
      \draw[accent,thick]
        plot[domain=0:12.6,samples=200,variable=\t]
          ({\r*exp(-0.3*\t)*cos(\t r)},{\r*exp(-0.3*\t)*(-sin(\t r))});}
    \fill[ink] (0,0) circle (1.5pt);
    \node[layergrey,anchor=north] at (0,-1.62)
      {(b) spiral stabil: $\lambda=-0.3\pm i$};
  \end{scope}
\end{tikzpicture}
\caption{Dua potret fase linear yang berbeda jenis. (a) Sistem
$x'=y$, $y'=x$: lintasan memenuhi $x^{2}-y^{2}=\text{konstan}$, hiperbola
dengan asimtot $y=\pm x$, dan garis putus-putus adalah arah eigen, satu masuk
dan satu keluar, sehingga titik asal adalah pelana dan tidak stabil. (b) Sistem $x'=-0{,}3x+y$, $y'=-x-0{,}3y$:
nilai eigen $-0{,}3\pm i$, sehingga lintasan berputar sambil menyusut dengan
faktor $e^{-0{,}3t}$ dan titik asal adalah spiral stabil. Empat lintasan di (b)
berbeda hanya pada jari-jari awal, karena sistemnya linear.}
\label{fig:potret-linear}
\end{figure}

Satu peringatan yang mahal. Bila $\tau=0$ dan $\Delta>0$, linearisasi memberi
pusat, tetapi sistem nonlinear di titik itu belum tentu pusat. Ambil
\[
x'=-y+x(x^{2}+y^{2}),\qquad y'=x+y(x^{2}+y^{2}).
\]
Jacobiannya di titik asal adalah matriks rotasi, nilai eigennya $\pm i$:
linearisasi berkata pusat. Tetapi dalam koordinat polar $r^{2}=x^{2}+y^{2}$,
perkalian silang memberi
\[
\frac{1}{2}\frac{d}{dt}r^{2}=xx'+yy'=r^{2}(x^{2}+y^{2})=r^{4},
\]
jadi $r'=r^{3}>0$ untuk $r>0$: lintasan berputar keluar, dan titik asal tidak
stabil. Linearisasi tidak salah; ia hanya tidak cukup ketika ada nilai eigen di
sumbu imajiner. Di situlah teorema tadi meminta titik setimbang hiperbolik.

Untuk kasus yang tidak hiperbolik, alat kedua dipakai: fungsi energi yang
menurun, yang kita karang sendiri.

\begin{definisi}[Fungsi Lyapunov]\index{fungsi Lyapunov}
Fungsi $V$ yang kontinu dan mempunyai turunan kontinu pada lingkungan $U$ dari
$\mathbf{x}^{*}$ disebut \emph{fungsi Lyapunov} bila $V(\mathbf{x}^{*})=0$ dan
$V(\mathbf{x})>0$ untuk $\mathbf{x}\in U\setminus\{\mathbf{x}^{*}\}$. Turunannya
sepanjang lintasan ditulis
\[
\dot V(\mathbf{x})=\nabla V(\mathbf{x})\cdot\mathbf{f}(\mathbf{x}).
\]
\end{definisi}

Bukit buatan itu adalah $V$, dan $\dot V$ kecuraman turunnya.

\begin{teorema}[Uji Lyapunov]
Misalkan $V$ fungsi Lyapunov di sekitar titik setimbang $\mathbf{x}^{*}$.
\begin{enumerate}
\item Bila $\dot V\le0$ pada $U$, maka $\mathbf{x}^{*}$ stabil.
\item Bila $\dot V<0$ pada $U\setminus\{\mathbf{x}^{*}\}$, maka
  $\mathbf{x}^{*}$ stabil asimtotik.
\end{enumerate}
\end{teorema}

\begin{proof}
Untuk pernyataan pertama, ambil $\varepsilon>0$ kecil sehingga bola
$\|\mathbf{x}-\mathbf{x}^{*}\|\le\varepsilon$ tinggal di dalam $U$. Di bola
itu $V$ positif dan bola itu kompak, jadi
$m=\min_{\|\mathbf{x}-\mathbf{x}^{*}\|=\varepsilon}V>0$. Pilih $\delta>0$
sehingga $V(\mathbf{x})<m$ untuk $\|\mathbf{x}-\mathbf{x}^{*}\|<\delta$;
mungkin karena $V$ kontinu dan $V(\mathbf{x}^{*})=0$. Lintasan dari bola $\delta$
memenuhi $V(\mathbf{x}(t))\le V(\mathbf{x}(0))<m$, sebab $\dot V\le0$,
sehingga ia tidak dapat menyeberangi bola $\varepsilon$; jadi stabil.

Untuk pernyataan kedua, $V(\mathbf{x}(t))$ menurun dan terbatas di bawah oleh
nol, jadi punya limit $L\ge0$. Andaikan $L>0$; lintasan itu tinggal di himpunan
kompak $\{V\ge L\}$, dan di sana $-\dot V$ punya minimum positif, sehingga $V$
harus turun melewati $L$ dalam waktu berhingga. Kontradiksi. Jadi $L=0$, dan
karena $V$ hanya nol di $\mathbf{x}^{*}$, lintasan menuju ke sana.
\end{proof}

Contohnya pendulum teredam, $\theta'=v$, $v'=-(g/L)\sin\theta-bv$ dengan
$b>0$, dan $V=v^{2}/2+(g/L)(1-\cos\theta)$, yaitu energi mekanik, positif
kecuali di titik setimbang. Turunannya
\[
\dot V=v\Bigl(-\frac{g}{L}\sin\theta-bv\Bigr)+\frac{g}{L}\sin\theta\,v=-bv^{2}\le0 .
\]
Jadi titik setimbang bawah stabil dan energi tidak pernah bertambah. Uji ini
tidak memberi kestabilan asimtotik, karena $\dot V$ nol di seluruh garis $v=0$;
versi LaSalle menambahkan bahwa satu-satunya himpunan invarian
di dalam $\{\dot V=0\}$ adalah titik setimbang, sehingga semua lintasan menuju
ke sana.

Linearisasi memberi jawaban yang sama. Pada $\theta=0$ Jacobiannya
$\begin{pmatrix}0&1\\-g/L&-b\end{pmatrix}$, dengan $\tau=-b$ dan
$\Delta=g/L>0$; untuk $b=0{,}5$ dan $g/L=9{,}80665$ akarnya $-0{,}25\pm3{,}121562\,i$,
yaitu spiral stabil. Pada $\theta=\pi$ muncul $+g/L$ di pojok kanan bawah,
sehingga $\Delta=-g/L<0$: pelana dengan nilai eigen $\pm3{,}131557$. Bandul yang
berdiri tegak memang pelana, bukan kiasan: ada satu arah tepat yang
mempertahankannya di atas, dan setiap gangguan membawanya pergi.

\kamus{stabil asimtotik}{asymptotically stable}{Tetap dekat setelah diganggu, lalu benar-benar tiba kembali ke titik setimbang.}

\kamus{fungsi Lyapunov}{Lyapunov function}{Bukit buatan yang dipakai sebagai penggaris: selama sistem menuruninya, ia tidak bisa lari jauh.}

\section{Euler, Runge--Kutta, dan kekakuan}
\tingkat{lanjut}

Hampir tidak ada persamaan yang dapat diselesaikan dengan rumus, jadi metode
numerik bukan pelengkap melainkan cara utama. Gagasannya kembali ke Bab 12: bila
kita tahu nilai dan turunan di satu titik, kita dapat menaksir nilai di sedikit
sebelah kanannya.

Metode Euler melakukan hal paling sederhana. Dari $y'=f(t,y)$ dan nilai awal,
\[
y_{n+1}=y_{n}+h\,f(t_{n},y_{n}),\qquad t_{n+1}=t_{n}+h .
\]
Ia tepat merupakan dua suku pertama deret Taylor di sekitar $t_{n}$, sehingga
galat satu langkahnya adalah suku $h^{2}y''(\xi)/2$, berorde $h^{2}$. Galat
lokal itu menumpuk menjadi galat global berorde $h$: menggandakan ketelitian
menuntut dua kali lipat jumlah langkah, bukan empat kali.

Angka untuk $y'=y$, $y(0)=1$, pada $t=1$ dengan $h=0{,}1$: Euler memberi
$(1{,}1)^{10}=2{,}5937424601$, sedangkan $e=2{,}7182818285$. Galatnya $0{,}1245$, yaitu
$4{,}58$ persen, dan galat satu langkah pertamanya $5{,}171\times10^{-3}$.

Metode Runge--Kutta orde empat memperbaiki ini dengan mengambil empat contoh
kemiringan di setiap langkah:
\[
y_{n+1}=y_{n}+\frac{h}{6}\bigl(k_{1}+2k_{2}+2k_{3}+k_{4}\bigr),
\]
dengan $k_{1}=f(t_{n},y_{n})$, $k_{2}=f(t_{n}+h/2,\,y_{n}+hk_{1}/2)$,
$k_{3}=f(t_{n}+h/2,\,y_{n}+hk_{2}/2)$, dan $k_{4}=f(t_{n}+h,\,y_{n}+hk_{3})$.
Galat lokalnya berorde $h^{5}$ dan galat globalnya berorde $h^{4}$: separuh
langkah berarti galat turun sekitar enam belas kali.

Tabel berikut berasal dari program di bagian lab, pada model logistik dengan
$y(0)=1/2$ dan $t=2$; nilai tepatnya $0{,}8807970780$.

\begin{center}
{\small
\begin{tabular}{lrrr}
\hline
$h$ & galat Euler & galat RK4 & rasio Euler \\
\hline
$0{,}400$ & $1{,}958\times10^{-2}$ & $1{,}245\times10^{-5}$ & --- \\
$0{,}200$ & $9{,}423\times10^{-3}$ & $7{,}237\times10^{-7}$ & $2{,}08$ \\
$0{,}100$ & $4{,}630\times10^{-3}$ & $4{,}359\times10^{-8}$ & $2{,}04$ \\
$0{,}050$ & $2{,}296\times10^{-3}$ & $2{,}674\times10^{-9}$ & $2{,}02$ \\
$0{,}025$ & $1{,}143\times10^{-3}$ & $1{,}656\times10^{-10}$ & $2{,}01$ \\
\hline
\end{tabular}}
\end{center}

Kolom terakhir adalah bukti orde pertama: setiap kali $h$ dibagi dua, galat
Euler turun hampir tepat dua kali; untuk RK4 rasionya sekitar $16$, sesuai orde
keempat. Gambar~\ref{fig:galat-numerik} menampilkan angka yang sama pada skala
logaritma kedua sumbu.

\begin{figure}[htbp]
\centering
\begin{tikzpicture}
\begin{axis}[width=0.80\linewidth,height=5.6cm,
  xlabel={langkah $h$},ylabel={$|y(2)-y_{h}(2)|$},
  xmode=log,ymode=log,xmin=0.02,xmax=0.5,ymin=1e-11,ymax=1e-1,
  grid=both,minor tick num=1,
  legend style={font=\scriptsize,at={(0.5,-0.22)},anchor=north,
    legend columns=2,draw=none,fill=none},
  tick label style={font=\scriptsize},label style={font=\small},
  axis line style={layergrey},tick style={layergrey}]
  \addplot[layerindigo,thick,mark=*,mark size=1.7pt] coordinates {
    (0.4,0.01958) (0.2,0.009423) (0.1,0.004630)
    (0.05,0.002296) (0.025,0.001143)};
  \addplot[accent,thick,mark=square*,mark size=1.6pt] coordinates {
    (0.4,0.00001245) (0.2,0.0000007237) (0.1,0.00000004359)
    (0.05,0.000000002674) (0.025,0.0000000001656)};
  \legend{Euler (kemiringan $1$), Runge--Kutta 4 (kemiringan $4$)}
\end{axis}
\end{tikzpicture}
\caption{Galat di $t=2$ untuk model logistik sebagai fungsi langkah $h$, pada
skala logaritma kedua sumbu. Garis Euler menurun dengan kemiringan satu: galat
berbanding lurus dengan $h$. Garis Runge--Kutta menurun dengan kemiringan empat:
galat berbanding dengan $h^{4}$. Pada $h=0{,}4$ kedua metode berjarak sekitar tiga
orde, dan jarak itu melebar ketika $h$ diperkecil. Angka pada gambar adalah
angka yang sama dengan tabel di teks.}
\label{fig:galat-numerik}
\end{figure}

Sekarang sifat yang membuat banyak orang tertipu: kekakuan. Sebuah sistem
disebut \emph{kaku} (\emph{stiff}) bila di dalamnya ada proses yang bergerak
jauh lebih cepat daripada proses yang kita minati. Ciri praktisnya: metode
eksplisit menuntut langkah sangat kecil untuk tetap \emph{stabil}, walaupun
ketelitian yang kita minta hanya menuntut langkah besar.

Contohnya kecil dan jujur: $y'=-50(y-\cos t)-\sin t$, $y(0)=1$, yang solusi
tepatnya $y=\cos t$. Di sini $y$ mengikuti $\cos t$ dengan lambat, tetapi setiap
penyimpangan diseret kembali dengan laju $50$ per detik. Euler stabil hanya bila
$|1-50h|<1$, yaitu $h<0{,}04$.

Hasil pada $t=3$, dengan nilai tepat $\cos3=-0{,}989992$: Euler dengan $h=0{,}02$
memberi galat $1{,}98\times10^{-4}$, dengan $h=0{,}039$ memberi
$7{,}97\times10^{-4}$, dan dengan $h=0{,}05$ memberi $-1{,}839\times10^{7}$. Satu
kenaikan dari $0{,}039$ ke $0{,}05$ mengubah jawaban dari hampir benar menjadi tidak
berarti. RK4 bertahan lebih lama, karena daerah stabilnya lebih lebar: pada
$h=0{,}05$ galatnya $1{,}11\times10^{-3}$, tetapi pada $h=0{,}1$ ia memberi
$-1{,}32\times10^{31}$; batas praktisnya $2{,}785/50=0{,}0557$.

Pelajaran yang dapat dibawa: galat dan kestabilan dua hal berbeda. Galat lokal
pada langkah pertama kecil dan tidak memperingatkan apa pun; yang meledak adalah
kestabilan, dan itu urusan nilai eigen.

Penawarnya metode implisit, yang paling sederhana Euler mundur:
\[
y_{n+1}=y_{n}+h\,f(t_{n+1},y_{n+1}),
\]
yang menuntut penyelesaian persamaan di setiap langkah, tetapi stabil untuk $h$
besar pada persamaan kaku; untuk contoh di atas, Euler mundur dengan $h=0{,}5$
memberi jawaban yang mendekati $\cos t$. Rujukan: Hairer dan Wanner, bab IV.

\begin{tantangan}
Linearisasi memberi pusat, tetapi sistem nonlinear tidak harus pusat. Susun satu
sistem dua dimensi yang titik asalnya \emph{bukan} pusat walaupun linearisasinya
pusat, tentukan sifat titik itu (stabil, stabil asimtotik, atau tidak stabil),
dan sertakan rumus eksplisit, nilai eigen Jacobian di titik asal, serta bukti
singkat klaimmu.

Petunjuk 1: linearisasi hanya melihat suku berderajat satu. Tambahkan suku
berderajat tiga yang simetris terhadap putaran, lalu tulis sistemnya dalam
koordinat polar.

Petunjuk 2: dalam koordinat polar, matriks rotasi menjadi $\theta'=1$, dan
gangguan radial menjadi $r'=cr^{3}$ untuk suatu konstanta $c$. Tanda $c$
menentukan segalanya, dan $V=r^{2}$ adalah fungsi Lyapunov yang bekerja untuk
salah satu tanda.

Kalau mentok, lihat kembali contoh $r'=r^{3}$ di bagian Kestabilan.
\end{tantangan}
\section{Contoh hitung}
\tingkat{menengah}

Empat hitungan berikut adalah tulang bab ini; semuanya dapat diperiksa ulang
dengan kalkulator atau program di bagian lab, dan semuanya pernah salah di tangan
saya sendiri.

\textbf{Model logistik: satu rumus, dua metode.}
Ambil $y'=y(1-y)$ dengan $y(0)=1/2$ dan tanyakan nilai di $t=2$. Rumus dari
bagian pemisahan variabel memberi $y(2)=1/(1+e^{-2})=0{,}8807970780$. Dengan
$h=0{,}1$ ada dua puluh langkah: Euler memberi $0{,}8854273272$, galat
$4{,}630\times10^{-3}$ atau $0{,}53$ persen, sedangkan Runge--Kutta dengan langkah
yang sama memberi $0{,}8807970344$, galat $4{,}359\times10^{-8}$. Biayanya empat kali
lipat, ketelitiannya sekitar seratus ribu kali.

Perhatikan juga arah galatnya. Euler di sini terlalu tinggi, karena
$y(1-y)$ cekung ke bawah pada selang itu sehingga garis lurus selalu berada di
atas kurva. Memeriksa tanda galat adalah cara murah untuk mengetahui apakah program kita
bekerja seperti seharusnya.

\textbf{Koefisien variabel: $y'+y/t=t^{2}$, $y(1)=0$.}
Di sini $p(t)=1/t$, sehingga
\[
\mu(t)=\exp\Bigl(\int\frac{dt}{t}\Bigr)=t
\]
pada $t>0$. Mengalikan persamaan dengan $t$ memberi
\[
t\,y'+y=t^{3}
\quad\Longleftrightarrow\quad
\bigl(t\,y\bigr)'=t^{3},
\]
jadi $t\,y=t^{4}/4+C$ dan
\[
y(t)=\frac{t^{3}}{4}+\frac{C}{t}.
\]
Syarat $y(1)=0$ memberi $C=-1/4$, sehingga
\[
y(t)=\frac{t^{3}}{4}-\frac{1}{4t}=\frac{t^{4}-1}{4t},
\qquad y(2)=\frac{15}{8}=1.875 .
\]
Pemeriksaan: $y'=3t^{2}/4+1/(4t^{2})$ dan $y/t=t^{2}/4-1/(4t^{2})$, jumlahnya
$t^{2}$, tepat ruas kanan. Integrator adaptif memberi $1{,}875000000000$.

Tanpa $\mu=t$, ruas kiri $y'+y/t$ tidak tampak seperti turunan apa pun; setelah
dikali $t$, ia menjadi turunan $ty$.

\textbf{Berapa besar ayunan sebelum rumus kecil meleset satu persen.}
Rasio periode tepat terhadap periode kecil adalah
\[
\frac{T(\theta_{0})}{T_{0}}
=\frac{2}{\pi}K\Bigl(\sin^{2}\frac{\theta_{0}}{2}\Bigr),
\]
dengan $K$ integral eliptik lengkap. Beberapa nilainya:

\begin{center}
{\small
\begin{tabular}{rrrr}
\hline
$\theta_{0}$ & $T/T_{0}$ & galat & keterangan \\
\hline
$5^{\circ}$ & $1{,}000476$ & $0{,}048$ persen & jam bandul \\
$10^{\circ}$ & $1{,}001907$ & $0{,}191$ persen & masih aman \\
$20^{\circ}$ & $1{,}007669$ & $0{,}767$ persen & mendekati batas \\
$22{,}81^{\circ}$ & $1{,}010000$ & $1{,}000$ persen & batas satu persen \\
$30^{\circ}$ & $1{,}017409$ & $1{,}741$ persen & perlu koreksi \\
$60^{\circ}$ & $1{,}073182$ & $7{,}318$ persen & jam meleset sejam \\
\hline
\end{tabular}}
\end{center}

Jadi rumus kecil itu layak dipakai sampai sekitar dua puluh derajat, dan
$22{,}81^{\circ}$ adalah ambang galat satu persen; galatnya tumbuh seperti kuadrat
amplitudo, bukan amplitudo.
Menggandakan amplitudo dari $10^{\circ}$ ke $20^{\circ}$ menaikkan galat dari
$0{,}19$ menjadi $0{,}77$ persen, yaitu empat kali, cocok dengan
$1+(\theta_{0}^{2}/16)$ sebagai dua suku pertama taksiran itu.

\kertas{
Uji sendiri dalam tiga langkah. (1) Kerjakan ulang contoh logistik dengan
$h=0{,}2$; engkau harus mendapat galat Euler $9{,}423\times10^{-3}$ dan galat RK4
$7{,}237\times10^{-7}$. (2) Turunkan sendiri rumus periode eliptik dari kekekalan
energi, lalu periksa bahwa $K$ pada parameter $0{,}25$ bernilai $1{,}6857503548$.
(3) Cari ambang $h$ untuk Euler pada $y'=-100y$; jawabannya $h<0{,}02$, dan pada
$h=0{,}02$ tepat, metode itu berayun antara $+1$ dan $-1$ tanpa meluruh maupun
meledak.}

\section{Lab komputasi}
\tingkat{lanjut}

Kode berikut menghitung semua angka di atas: galat Euler dan RK4 pada model
logistik, periode pendulum dengan dua jalan berbeda, dan ambang kekakuan. Ia
hanya memakai NumPy dan SciPy.

\begin{lab}
\begin{verbatim}
import numpy as np
from scipy import integrate, special

def euler(f, y0, T, h):
    y, t, n = y0, 0.0, int(round(T / h))
    for _ in range(n):
        y += h * f(t, y); t += h
    return y

def rk4(f, y0, T, h):
    y, t, n = y0, 0.0, int(round(T / h))
    for _ in range(n):
        k1 = f(t, y); k2 = f(t + h/2, y + h*k1/2)
        k3 = f(t + h/2, y + h*k2/2); k4 = f(t + h, y + h*k3)
        y += h * (k1 + 2*k2 + 2*k3 + k4) / 6; t += h
    return y

logi = lambda t, y: y * (1 - y)
ex = lambda t: 1 / (1 + np.exp(-t))
print("logistik y(2) eksak %.10f" % ex(2.0))
for h in (0.4, 0.2, 0.1, 0.05, 0.025):
    ee = abs(euler(logi, 0.5, 2, h) - ex(2))
    er = abs(rk4(logi, 0.5, 2, h) - ex(2))
    print("h=%.3f euler %.3e rk4 %.3e" % (h, ee, er))

g, L = 9.80665, 1.0
T0 = 2 * np.pi * np.sqrt(L / g)
th0 = np.radians(60.0)
k = np.sin(th0 / 2)
Tq = 4 * np.sqrt(L / (2*g)) * integrate.quad(
        lambda th: 1/np.sqrt(np.cos(th) - np.cos(th0)), 0, th0, limit=200)[0]
Te = 4 * np.sqrt(L / g) * special.ellipk(k**2)
print("T0 %.6f T60kuad %.6f T60eliptik %.6f" % (T0, Tq, Te))
print("rasio %.6f persen %.4f" % (Te / T0, 100 * (Te / T0 - 1)))

stiff = lambda t, y: -50 * (y - np.cos(t)) - np.sin(t)
for h in (0.02, 0.039, 0.05):
    print("kaku euler h=%.3f y(3)=%.3e" % (h, euler(stiff, 1.0, 3, h)))
for h in (0.05, 0.1):
    print("kaku rk4 h=%.3f y(3)=%.3e" % (h, rk4(stiff, 1.0, 3, h)))
print("eksak cos(3) %.10f" % np.cos(3))
\end{verbatim}
\end{lab}

Keluaran sebenarnya, apa adanya:

\begin{verbatim}
logistik y(2) eksak 0.8807970780
h=0.400 euler 1.958e-02 rk4 1.245e-05
h=0.200 euler 9.423e-03 rk4 7.237e-07
h=0.100 euler 4.630e-03 rk4 4.359e-08
h=0.050 euler 2.296e-03 rk4 2.674e-09
h=0.025 euler 1.143e-03 rk4 1.656e-10
T0 2.006409 T60kuad 2.153242 T60eliptik 2.153242
rasio 1.073182 persen 7.3182
kaku euler h=0.020 y(3)=-9.902e-01
kaku euler h=0.039 y(3)=-9.908e-01
kaku euler h=0.050 y(3)=-1.839e+07
kaku rk4 h=0.050 y(3)=-9.883e-01
kaku rk4 h=0.100 y(3)=-1.317e+31
eksak cos(3) -0.9899924966
\end{verbatim}

Tiga hal patut dibaca. Pertama, rasio galat Euler $2{,}08$, $2{,}04$, $2{,}02$, $2{,}01$
dan RK4 menuju $16=2^{4}$: orde pertama dan keempat terlihat sebagai angka. Kedua,
dua jalan untuk periode pendulum bertemu sampai enam digit, padahal kuadratur dan
integral eliptik tidak memakai rumus yang sama. Ketiga, baris kaku: $h=0{,}039$
masih wajar, $h=0{,}05$ tidak berarti, dan tidak satu pun baris memperingatkan
kita.

Program itu tidak memeriksa syarat teorema: apakah $f$ memenuhi Lipschitz,
apakah titik setimbangnya hiperbolik. Itu tetap pekerjaan kita.

\begin{untukdosen}
Urutan mengajar yang saya pakai: mulai dari jam bandul dan selisih $7{,}32$
persen, lalu definisi solusi beserta contoh ledakan $1/(1-t)$, baru pemisahan
variabel dan faktor integrasi. Faktor integrasi sebaiknya dibuktikan penuh di
papan, karena buktinya hanya memakai aturan hasil kali dan teorema dasar
kalkulus, dua alat yang sudah mereka punya sejak Bab 13. Picard--Lindelof cukup
sebagai sketsa dengan satu iterasi dijalankan; yang harus tinggal adalah arti
syarat Lipschitz, bukan teknik titik tetapnya. Bidang fase sebelum nilai eigen:
gambar dulu, hitung kemudian.

Dua miskonsepsi yang hampir selalu muncul. Pertama, mahasiswa menganggap
"solusi" berarti rumus tertutup, lalu menyimpulkan persamaan tanpa rumus tidak
punya solusi; contoh $y'=\sqrt{|y|}$ membantu, karena di sana ada dua solusi dan
tidak ada rumus tunggal. Kedua, mereka memperlakukan galat kecil sebagai bukti
kebenaran, dan terkejut ketika Euler dengan $h=0{,}05$ memberi $-1{,}8\times10^{7}$.
Tunjukkan baris keluaran itu, jangan hanya diperingatkan.

Satu aktivitas papan tulis: gambar bidang fase pendulum kosong, lalu minta
mahasiswa menandai di mana kecepatan nol, di mana percepatan nol, dan di mana
energi tertinggi. Kelas akan berdebat tentang daerah di antara dua pelana, dan
perdebatan itu persis isi teorema. Tanda kesiapan lanjut: mahasiswa dapat
menebak jenis titik setimbang dari jejak dan determinan tanpa menghitung akar.

\end{untukdosen}

\section{Ringkasan}
\begin{ringkasan}
\begin{itemize}
\item Persamaan diferensial menyamakan laju perubahan dengan keadaan; solusinya
fungsi pada sebuah selang, dan selang itu bagian dari jawaban.
\item Dua keluarga terpecahkan dengan tangan: pemisahan variabel dan persamaan
linear orde satu dengan faktor integrasi $\mu=\exp(\int p\,dt)$.
\item Picard--Lindel\"of menjamin solusi tunggal pada selang pendek bila $f$
kontinu dan memenuhi syarat Lipschitz terhadap $y$; tanpa syarat itu, satu
keadaan awal dapat melahirkan dua riwayat.
\item Bidang fase mengganti waktu dengan gambar: titik setimbang, lintasan
tertutup, separatriks, dan lintasan terbuka terbaca tanpa rumus solusi.
\item Kestabilan lokal dibaca dari nilai eigen Jacobian, kecuali bila ada nilai
eigen di sumbu imajiner; di situ linearisasi tidak memutuskan, dan fungsi
Lyapunov dipakai.
\item Euler berorde satu secara global dan RK4 berorde empat; keduanya eksplisit
dan karena itu gagal pada sistem kaku, tempat kestabilan, bukan ketelitian,
yang menentukan besar langkah.
\end{itemize}
\end{ringkasan}

\section{Latihan}
\begin{latihan}
\begin{enumerate}
\item Selesaikan $y'=2ty$, $y(0)=1$, lalu hitung $y(1)$ dan bandingkan dengan
  $e$. Hitung pula Euler dengan $h=0{,}1$ dan sebutkan galatnya. (Hasil saya:
  $2{,}33463336$, galat $3{,}84\times10^{-1}$.)
\item Buktikan bahwa $y'+p(t)y=q(t)$ dengan $p,q$ kontinu punya paling banyak
  satu solusi, tanpa rumus eksplisit: selisih dua solusi memenuhi
  $w'=-p(t)w$ dan $w(t_{0})=0$.
\item Benar, salah, atau belum cukup: ``karena linearisasi di titik asal memberi
  nilai eigen $\pm i$, titik asal adalah pusat yang stabil''. Beri satu contoh
  pendukung atau satu penyangkal.
\item Buktikan $E=v^{2}/2+\omega_{0}^{2}\theta^{2}/2$ konstan pada osilator
  harmonik, lalu jelaskan mengapa lintasannya elips dan periodenya
  $2\pi/\omega_{0}$.
\item Rangkaian RC dengan $R=2{,}2\,\mathrm{k}\Omega$, $C=10\,\mu\mathrm{F}$,
  $V_{s}=12$ volt, mulai dari nol. Hitung $\tau$, tegangan pada $t=\tau$, dan
  waktu sampai $6$ volt. ($22$ milidetik; $7{,}585447$ volt; $15{,}249$ milidetik.)
\item Tunjukkan bahwa $v=\pm2\cos(\theta/2)$ memenuhi $E=2$, lalu jelaskan
  mengapa kurva itu menghubungkan $(-\pi,0)$ dengan $(\pi,0)$ dan disebut
  separatriks.
\item Bandul $L=0{,}25$ meter dilepas dari $30^{\circ}$. Hitung periode kecil dan
  periode tepatnya, lalu sebutkan selisih persennya. ($1{,}003205$ detik;
  $1{,}020669$ detik; $1{,}74$ persen.)
\item Untuk $y'=-100y$ dengan $y(0)=1$, tentukan ambang $h$ yang membuat Euler
  eksplisit tetap stabil, lalu jelaskan apa yang terjadi tepat pada ambang itu.
  Periksa dengan program.
\item Buktikan bahwa persamaan otonom skalar $y'=f(y)$ dengan $f$ kontinu dan
  memenuhi Lipschitz tidak dapat mempunyai solusi periodik tidak konstan.
\end{enumerate}

\textbf{Petunjuk tiga butir tersulit.} Butir 3: linearisasi hanya melihat suku
berderajat satu; tambahkan suku kubik simetris seperti pada contoh $r'=r^{3}$
dan bandingkan dengan $r'=-r^{3}$. Butir 6: pakai
$\cos\theta=2\cos^{2}(\theta/2)-1$ dan periksa nilai $v$ di $\theta=\pm\pi$.
Butir 9: andaikan $y$ periodik tidak konstan; di titik maksimum $y'=0$, jadi
$f(y_{\max})=0$, dan ketunggalan memaksa lintasan itu konstan.
\end{latihan}

\begin{untukpeneliti}
Yang benar-benar terbukti di bab ini, dengan bukti yang ditulis: rumus faktor
integrasi beserta ketunggalannya, sketsa lengkap Picard--Lindelof dengan
batas galat iterasi, ketunggalan solusi linear lewat $w'=-pw$, uji Lyapunov
untuk kedua pernyataannya, dan seluruh identitas energi pada osilator serta
pendulum. Yang saya kutip tanpa bukti: teorema Hartman--Grobman (Hartman,
\emph{Ordinary Differential Equations}, bab 9; Perko, \emph{Differential
Equations and Dynamical Systems}, bagian 2.8), teorema Peano dan
Arzel\`a--Ascoli, serta prinsip invarians LaSalle. 

Yang dipopulerkan secara berlebihan adalah ungkapan "nilai eigen menentukan
segalanya". Yang benar: nilai eigen menentukan perilaku lokal hanya bila tidak
ada yang bagian realnya nol; di sumbu imajiner suku berderajat tiga dapat
mengubah pusat menjadi spiral, seperti contoh $r'=r^{3}$ di sini. Untuk sistem
berdimensi tinggi, "stabil asimtotik" lokal bahkan tidak berarti aman dari
gangguan besar, karena daerah tariknya dapat sangat kecil.

Pertanyaan terbuka: seberapa jauh sifat kualitatif dapat disertifikasi secara
numerik? Galat kecil pada satu lintasan tidak memberi tahu apa pun tentang
kestabilan; Euler dengan $h=0{,}039$ benar sampai empat digit, sedangkan $h=0{,}05$
memberi $10^{7}$. Untuk sistem yang hanya diketahui lewat sampel medan, tidak ada
cara umum memastikan tidak ada daerah kaku yang tersembunyi di antara dua titik
sampel.

Batas keberlakuan: semua fungsi di sini real dan cukup mulus; semua angka dari
Python 3.14.4, NumPy 2.5.2, dan SciPy 1.18.0, dengan kuadratur bertoleransi
$10^{-13}$ kecuali disebut lain; $g=9{,}80665\ \mathrm{m/s^{2}}$ dan gesekan
diabaikan pada hitungan periode.
\end{untukpeneliti}

\peta{Bab 16 pada Boyce dan DiPrima, \emph{Elementary Differential Equations
and Boundary Value Problems}, bab 1 (pengantar), 2 (persamaan orde satu), 3
(persamaan linear orde dua), 4 (sistem orde satu dan bidang fase), dan 9
(persamaan nonlinear dan kestabilan); pada Strogatz, \emph{Nonlinear Dynamics
and Chaos}, bab 2 (aliran pada garis), 5 (sistem dua dimensi), 6 (bidang fase),
7 (osilator), dan 8 (bifurkasi); pada Arnold, \emph{Ordinary Differential
Equations}, bab 1--4 untuk keberadaan, ketunggalan, dan klasifikasi titik
setimbang; pada Coddington dan Levinson, bab 1--2 untuk Picard--Lindelof dan
Peano; pada Hartman, bab 9, dan Perko, bagian 2.8, untuk Hartman--Grobman;
pada Hairer dan Wanner, \emph{Solving Ordinary Differential Equations II}, bab
IV, untuk kekakuan dan metode implisit. Dasar hitungannya bab \emph{Turunan} (turunan parsial dan Jacobian), bab
\emph{Deret Fungsi dan Pertukaran Limit} (limit seragam, yang membuat iterasi
Picard konvergen), dan bab \emph{Integral Lipat} (medan vektor) pada edisi
pertama buku ini; nilai eigen dari Bab 9 dipakai langsung tanpa diulang, dan
analisis galat floating-point dari Bab 22 menjadi lanjutan bagian numerik bab
ini.}
